% s-211-p211-m002; original PDF 211, 212
\Aphorism{s-211-p211-m002}{305}
In unserem Jahrhundert sehe ich zwei grosse Künstler, von denen man ohne sittliches Pathos gar nicht reden kann, Künstler, die schlechthin unbegriffen bleiben, solange man sie als rein artistische oder intellectuelle Begabungen begreifen und nicht auch ihrem Menschlichen gerecht werden will: Richard Wagner und Friedrich Hebbel. Ich denke hier nicht an Stoffliches und Technisches, worin sich beide Männer berühren, nicht dass Beide für ihre Kunst das Reich der dialectischen Leidenschaft und der erotischen Überbewusstheit entdeckt haben, dass Jeder von ihnen die Grenzen seiner Sprache erweitert und die Concision des Ausdrucks verdoppelt hat, noch weniger natürlich daran, dass Beide einen Nibelungencyclus gedichtet, dramaturgische Theorien aufgestellt und sich mit inferioren Zeitgenossen erbärmlich geplagt haben: sondern ich denke an das Gemeinsame des Wollens, das hier einmal zwei gewaltigen Könnern ersten Ranges anhaftet, während es sonst gern die vom zweiten Range zu einer Genossenschaft verbindet, der die führenden Geister fernbleiben. Hebbel und Wagner sind, so nah oder fern sie sich sonst stehen mögen, Eins in der \Sperr{ethischen Auffassung} des künstlerischen Berufes, im feierlichen Ernst ihres Glaubens an sich selbst, an ihre Sendung und Aufgabe, in der unbedingten Abkehr von einem entweihten und verwahrlosten Kunstzeitalter. Diese Art Mensch (unter den Künstlern zweiten Ranges, wie gesagt, weniger selten als unter den ganz Grossen) ist mir immer ehrwürdig erschienen durch die innere Nothwendigkeit im Thun und Schaffen; sie ist am fernsten von der Kunst als Industriezweig, fern von der Kunst als Stubenhockerei und Schwärmerei, fern aber auch von der „Kunst für Künstler“, von jeder wenn auch noch so verfeinerten und verklärten hedonistischen Werthung, vom Sybaritismus und von der Gourmandise. Wagner und Hebbel glaubten an die lebendige Kunst als an ein sociales Ideal, einen Mittelpunkt, um den sich von selbst die anderen Culturgewalten schaaren müssten, und Beide sahen die Anknüpfungsstelle für ihre Hoffnungen da, wo der Augenschein sie nicht zu suchen lehrte: im Theater. Beide erkannten dem Drama die höchste Mission zu, an die zu glauben man wider die modernen Verfallszustände empfinden muss; Beide hofften in verwandter Weise vom Volke und vom deutschen Wesen, was angesichts der modernen Zustände kein Anderer zu hoffen wagte. Beide dienten der Kunst und kannten nur Eine unverzeihliche Sünde: Verrath an der Kunst. Es sind die beiden Evangelisten der Kunst im neunzehnten Jahrhundert: wo bleiben die übrigen?
\par

